\chapter{Stand der Technik}
\label{ch:StandDerTechnik}

\section{Prognosen mit fehlenden und auffälligen Werten}

Die Anomalieerkennung mit Vorhersagen setzt voraus, dass ein Modell das
Normalverhalten ausreichend lernt. In unbeschrifteten Trainingsdaten können
jedoch bereits unbekannte Anomalien enthalten sein. RobustTSF zeigt für mehrere
dort definierte Zielanomalien Vorteile des absoluten gegenüber dem quadratischen
Fehler \cite{cheng2024robusttsf}. DCFF-MTAD berichtet in drei Vergleichstests
Verbesserungen durch die Huber-Fehlerfunktion, die große Abweichungen weniger
stark gewichtet als MSE \cite{xu2023dcffmtad}. Daraus folgt keine allgemein
beste Fehlerfunktion, wohl aber die Notwendigkeit, Vorhersagegüte und
Anomalieerkennung getrennt zu bewerten.

Fehlende Werte tragen außerdem selbst Information. GRU-D verwendet deshalb
unter anderem eine Beobachtungsmaske und zeigt, dass ein rekurrentes Modell
Muster fehlender Messungen nutzen kann \cite{che2018grud}. Eine solche Maske
ist deutlich einfacher als die vollständige GRU-D-Struktur und eignet sich
daher auch für ein kleines Modell.

\section{Kleine rekurrente Modelle}

Das LSTM wurde entwickelt, um längerfristige Abhängigkeiten in Sequenzen zu
lernen \cite{hochreiter1997long}. Für die nur zwei Wochen langen Sensorreihen
ist ein großes Netz jedoch nicht automatisch geeigneter. Dropout setzt während
des Trainings zufällig Aktivierungen auf null und wurde als einfache Maßnahme
gegen Überanpassung eingeführt \cite{srivastava2014dropout}. Das spricht für
ein kleines Netz mit mäßigem Dropout, legt aber keine bestimmte Anzahl von
Einheiten fest.

\section{Bewertung der Anomalieerkennung}

Eine hohe Vorhersagegenauigkeit allein garantiert keine gute
Anomalieerkennung. AFMF zeigt, dass unter anderem Skalierungsunterschiede und
Anomalien im Eingabefenster die Unterscheidung zwischen normalem und
auffälligem Verhalten erschweren können \cite{shen2024afmf}.
Goswami et al. untersuchen künstlich erzeugte Anomalien als
Vergleichsmöglichkeit ohne reale Kennzeichnungen und finden einen deutlichen
Zusammenhang mit der Bewertung auf beschrifteten Daten
\cite{goswami2023unsupervised}. Reale Anomalien können von den künstlichen
Typen abweichen; solche Tests ersetzen deshalb keine fachliche Prüfung.

mTSBench vergleicht 24 Verfahren auf 344 mehrdimensionalen Zeitreihen aus 19
Datensätzen. Kein Verfahren dominiert auf allen Datensätzen
\cite{zhou2026mtsbench}. Ein Modell muss daher auf den konkreten Daten geprüft
werden; gute Ergebnisse auf einem fremden Datensatz lassen sich nicht einfach
übertragen.

\section{Prognose und Rekonstruktion}

MTAD-GAT optimiert Prognose und Rekonstruktion gemeinsam
\cite{zhao2020mtadgat}. Im Vergleich innerhalb dieser Studie ist die gemeinsame
Nutzung beider Aufgaben auf drei Datensätzen besser als jeweils nur eine
Aufgabe. Das garantiert keinen Gewinn für Bad Nauheim, zeigt aber, dass
Prognose- und Rekonstruktionsfehler unterschiedliche Hinweise auf Anomalien
liefern können.
